% IEEE conference version for GTSD 2026 proceedings
\documentclass[conference]{IEEEtran}
\usepackage{cite}
\usepackage{amsmath,amssymb,amsfonts}
\usepackage{graphicx}
\graphicspath{{figures/}}
\usepackage{textcomp}
\usepackage{booktabs}
\usepackage{url}
\def\BibTeX{{\rm B\kern-.05em{\sc i\kern-.025em b}\kern-.08em
    T\kern-.1667em\lower.7ex\hbox{E}\kern-.125emX}}

\begin{document}
\bstctlcite{IEEEexample:BSTcontrol}

% Initialize section counter to 3 so this section compiles cleanly as Section IV
\setcounter{section}{3}
\section{Experimental Results and Analysis}
\label{sec:experiments}

% Initialize counters to maintain continuity with Section III
\setcounter{subsection}{0}
\setcounter{equation}{20}  % Continues from Section III-E
\setcounter{figure}{6}    % Fig. 7 onwards
\setcounter{table}{1}     % Table II onwards

This section establishes the empirical evaluation framework to assess the image quality, timing throughput, power dissipation, and hardware resource efficiency of the proposed FPGA accelerator. The experimental analysis is organized into the following structured subsections:

% =========================================================================
% IV-A. EXPERIMENTAL SETUP AND BENCHMARK DATASETS (OUTLINE)
% =========================================================================
\subsection{Experimental Setup and Benchmark Datasets}
\label{sec:exp_setup}

\subsubsection{Hardware Platform and Environment Setup}
\begin{itemize}
    \item \textbf{Target SoC:} Xilinx PYNQ-Z2 board featuring the Zynq-7020 SoC (\texttt{xc7z020clg400-1}).
    \item \textbf{Processing System (PS):} Dual-core ARM Cortex-A9 @ 650~MHz, 512~MB DDR3, running embedded PYNQ Linux (v2.6). Responsible for file I/O, initial bicubic scaling, DMA memory buffer management, and final image stitching.
    \item \textbf{Programmable Logic (PL):} Artix-7 FPGA fabric running at a unified clock of 100~MHz. Connected to PS via 64-bit AXI4-Stream DMA over the High-Performance (HP) port.
    \item \textbf{Synthesis and Implementation Tools:} AMD Xilinx Vivado Design Suite 2020.2; power analysis using Vivado Power Sign-off Analyzer under post-route physical constraints.
\end{itemize}

\subsubsection{Benchmark Clinical Datasets}
\begin{itemize}
    \item \textbf{NIH ChestX-ray14 \cite{wang2017chestxray}:} Independent test partition of adult chest radiographs ($1024 \times 1024$ pixels) covering diverse thoracic pathologies.
    \item \textbf{Guangzhou Pediatric Pneumonia \cite{kermany2018identifying}:} Pediatric chest radiographs of varying clinical resolutions (ranging from $1000 \times 800$ up to $1638 \times 1236$ pixels) to evaluate arbitrary-dimension tiling robustness.
    \item \textbf{Calibration Subset \cite{chowdhury2020can}:} Used exclusively for initial fixed-point dynamic range verification and activation quantization calibration.
\end{itemize}

\subsubsection{Comparative Evaluation Baselines}
\begin{itemize}
    \item \textbf{Bicubic Upsampling ($2\times$):} Standard software baseline widely used in edge medical systems.
    \item \textbf{FP32 Compact SRCNN:} Software floating-point baseline executed on CPU/GPU.
    \item \textbf{Fixed-Point Emulation (PTQ Q7 Software):} Bit-accurate Python emulation modeling the exact integer arithmetic of the RTL core.
    \item \textbf{Proposed FPGA Accelerator:} Hardware execution on the Xilinx Zynq-7020 PL fabric.
\end{itemize}

% =========================================================================
% IV-B. IMAGE QUALITY AND QUANTIZATION FIDELITY (OUTLINE)
% =========================================================================
\subsection{Image Quality and Quantization Fidelity Evaluation}
\label{sec:quality_eval}

\subsubsection{Quantitative Metrics Protocol}
To evaluate image restoration quality and quantization-induced precision loss, the following objective metrics will be reported:
\begin{itemize}
    \item Peak Signal-to-Noise Ratio (PSNR in dB) and Structural Similarity Index (SSIM).
    \item Quantization Loss: $\Delta\text{PSNR} = \text{PSNR}_{\text{FPGA}} - \text{PSNR}_{\text{FP32}}$ and $\Delta\text{SSIM} = \text{SSIM}_{\text{FPGA}} - \text{SSIM}_{\text{FP32}}$.
    \item Performance gain over standard bicubic interpolation ($\text{PSNR}_{\text{FPGA}} - \text{PSNR}_{\text{Bicubic}}$).
\end{itemize}

\begin{table*}[t]
\caption{Quantitative Image Quality and Quantization Degradation Framework Across Clinical Datasets [Outline Skeleton]}
\label{tab:image_quality_benchmark}
\centering
\begin{tabular}{lc|cc|cc|cc|cc}
\hline
\textbf{Clinical Dataset} & \textbf{Evaluated} & \multicolumn{2}{c|}{\textbf{Bicubic Baseline}} & \multicolumn{2}{c|}{\textbf{Compact SRCNN (FP32)}} & \multicolumn{2}{c|}{\textbf{Proposed FPGA (PTQ Q7)}} & \multicolumn{2}{c}{\textbf{Quantization Loss ($\Delta$)}} \\
\textbf{Benchmark} & \textbf{Samples} & \textbf{PSNR (dB)} & \textbf{SSIM} & \textbf{PSNR (dB)} & \textbf{SSIM} & \textbf{PSNR (dB)} & \textbf{SSIM} & \textbf{$\Delta$PSNR (dB)} & \textbf{$\Delta$SSIM} \\
\hline
NIH ChestX-ray14 \cite{wang2017chestxray} & [Pending] & [TBD] & [TBD] & [TBD] & [TBD] & [TBD] & [TBD] & [TBD] & [TBD] \\
Guangzhou Pediatric \cite{kermany2018identifying} & [Pending] & [TBD] & [TBD] & [TBD] & [TBD] & [TBD] & [TBD] & [TBD] & [TBD] \\
\hline
\textbf{Overall Clinical Benchmark} & \textbf{[Total]} & \textbf{[TBD]} & \textbf{[TBD]} & \textbf{[TBD]} & \textbf{[TBD]} & \textbf{[TBD]} & \textbf{[TBD]} & \textbf{[TBD]} & \textbf{[TBD]} \\
\hline
\end{tabular}
\end{table*}

\subsubsection{Qualitative Analysis and Diagnostic Safety Verification}
\begin{itemize}
    \item \textbf{Visual Comparison (Planned Fig.~7):} Zoomed-in Regions of Interest (ROIs) on critical anatomical structures (rib trabeculae, pulmonary vessels, pleural boundaries) comparing Bicubic, FP32, and FPGA outputs.
    \item \textbf{Hallucination-Avoidance Verification:} Demonstration that deterministic mapping with bounded fixed-point dynamic range avoids the synthetic structural artifacts common in GAN-based and diffusion-based super-resolution models.
    \item \textbf{Residual Error Check:} Verification of pixel-level error bounds ($\text{MAE} < 1\text{ LSB}$) and confirmation of zero overflow/clipping events during physical execution.
\end{itemize}

% =========================================================================
% IV-C. BOUNDARY-ARTIFACT ELIMINATION (OUTLINE)
% =========================================================================
\subsection{Boundary-Artifact Elimination and Overlap-Tiling Verification}
\label{sec:tiling_eval}

\begin{itemize}
    \item \textbf{Evaluation Objective:} Empirically validate Algorithm~1 (Hardware-Aware Overlap-Tiling) in preventing edge truncation artifacts across tile boundaries.
    \item \textbf{Comparative Setup:}
    \begin{itemize}
        \item \textit{Non-overlapping Tiling ($S=128, M=0$):} Observation of boundary seam discontinuities and localized PSNR drops along patch seams due to receptive field truncation ($M_{rf} = 6\text{ px}$).
        \item \textit{Proposed Overlap-Tiling ($S=112, M=8$):} Elimination of boundary seams through 8-pixel margin exclusion and edge-replication padding.
    \end{itemize}
    \item \textbf{Stitching Overhead Analysis:} Assessment of PS canvas reconstruction latency (approximately 1.54\% of full-frame time for $1024 \times 1024$ frames).
\end{itemize}

% =========================================================================
% IV-D. END-TO-END SYSTEM LATENCY AND PROFILING (OUTLINE)
% =========================================================================
\subsection{End-to-End System Latency and Execution Profiling}
\label{sec:latency_eval}

\subsubsection{Latency Decomposition Framework}
Measurement and breakdown of execution time across the complete heterogeneous pipeline:
\begin{itemize}
    \item $T_{\text{bicubic}}$: Host-side bicubic pre-interpolation on the ARM Cortex-A9.
    \item $\sum T_{\text{DMA+PL}}$: DMA packet streaming and pipelined PL compute across all image patches.
    \item $T_{\text{stitch}}$: Canvas reconstruction, margin exclusion, and memory assembly on the host.
\end{itemize}

\begin{table*}[t]
\caption{End-to-End Execution Latency and Throughput Breakdown Framework [Outline Skeleton]}
\label{tab:latency_breakdown}
\centering
\begin{tabular}{lcccccc}
\hline
\textbf{Resolution} & \textbf{Tiling Scheme} & \textbf{Tile Count} & \textbf{Bicubic Pre-Proc.} & \textbf{AXI DMA + PL Compute} & \textbf{Canvas Stitching} & \textbf{Total Frame Time} \\
\textbf{($H \times W$)} & \textbf{\& Pitch} & \textbf{($N_p$)} & \textbf{(PS ARM Cortex-A9)} & \textbf{(64-bit HP Bus + PL)} & \textbf{(PS ARM Cortex-A9)} & \textbf{(Frame Rate)} \\
\hline
$1024 \times 1024$ & Non-overlapping ($S=128$) & 64 & [TBD] & [TBD] & [TBD] & [TBD FPS] \\
$1024 \times 1024$ & Overlap-Tiling ($S=112$) & 100 & [TBD] & [TBD] & [TBD] & [TBD FPS] \\
$1638 \times 1236$ & Overlap-Tiling ($S=112$) & 180 & [TBD] & [TBD] & [TBD] & [TBD FPS] \\
\hline
\end{tabular}
\end{table*}

\subsubsection{Analysis of Memory Hierarchy and Scaling Overhead}
\begin{itemize}
    \item Detailed scaling analysis of DMA transfers and PL execution vs. tile count.
    \item Investigation of ARM Cortex-A9 memory cache effects during non-contiguous 2D strided array stitching on large image canvases ($>1.5\text{ Mpixels}$).
    \item End-to-end hardware acceleration speedup relative to pure ARM CPU software execution.
\end{itemize}

% =========================================================================
% IV-E. COMPARISON WITH STATE-OF-THE-ART (OUTLINE)
% =========================================================================
\subsection{Hardware Resource Utilization and SOTA Comparison}
\label{sec:sota_comparison}

\subsubsection{Comparison Protocol}
To benchmark the proposed accelerator against recent literature, this subsection will extend the analysis initiated in Section~II (Table~I), comparing physical implementation metrics against state-of-the-art FPGA super-resolution designs.

\begin{table*}[t]
\caption{Comprehensive Comparison with State-of-the-Art FPGA Super-Resolution Accelerators [Outline Skeleton]}
\label{tab:sota_extended}
\centering
\begin{tabular}{lcccccc}
\hline
\textbf{Design Metric} & \textbf{Kim \textit{et al.} \cite{kim2019real}} & \textbf{Mao \textit{et al.} \cite{mao2020fdna}} & \textbf{Liu \textit{et al.} \cite{liu2024high}} & \textbf{Chang \textit{et al.} \cite{chang2024hdsuper}} & \textbf{Huang \textit{et al.} \cite{huang2026uasr}} & \textbf{This Work} \\
\hline
Publication Year & 2019 & 2020 & 2024 & 2024 & 2026 & \textbf{2026} \\
Target FPGA Device & Xilinx Virtex-7 & Xilinx Kintex-7 & Xilinx ZCU102 & Xilinx XC7Z035 & Xilinx ZCU102 & \textbf{Xilinx Zynq-7020} \\
Platform Classification & High-End FPGA & Mid-Tier FPGA & High-End MPSoC & Mid-Tier SoC & High-End MPSoC & \textbf{Low-Cost Edge SoC} \\
Clock Frequency ($F_{clk}$) & 150~MHz & 200~MHz & 200~MHz & 142~MHz & 250~MHz & \textbf{100~MHz} \\
DSP48 Utilization & 740 (26.4\%) & 1,512 (98.2\%) & 1,210 (48.0\%) & 704 (78.2\%) & 181 (7.2\%) & \textbf{120 (54.55\%)} \\
BRAM36K Blocks & 432 & 235 & 412 & 312 & 142 & \textbf{3.50 (2.50\%)} \\
LUT Utilization & 142,500 & 82,410 & 118,400 & 94,200 & 42,100 & \textbf{8,615 (16.20\%)} \\
Total SoC Power & 12.80~W & 5.40~W & 16.85~W & 2.08~W & 2.88~W & \textbf{1.438~W} \\
PL Dynamic Power & --- & 3.82~W & 14.20~W & 1.34~W & 1.95~W & \textbf{0.047~W} \\
Energy Efficiency & 4.69 FPS/W & 11.61 FPS/W & 2.12 FPS/W & 38.94 FPS/W & 28.82 FPS/W & \textbf{[To Report]} \\
Throughput & 60.0~fps (4K) & 62.7~fps (UHD) & 35.8~fps (FHD) & 81.0~fps (FHD) & 83.0~fps (FHD) & \textbf{157.9~p/s} \\
Quantization Scheme & Fixed 16-bit & Fixed 13-bit & Fixed 8-bit & Fixed 16-bit & Fixed 16-bit & \textbf{PTQ Q7 (INT8)} \\
Target Application & Natural Video & Natural Images & Natural Images & Natural Images & Natural Images & \textbf{Medical Radiographs} \\
Clinical Validation & No & No & No & No & No & \textbf{Yes (Clinical CXR)} \\
\hline
\end{tabular}
\end{table*}

\subsubsection{Key Architectural Comparative Highlights}
\begin{itemize}
    \item \textbf{DSP and Memory Footprint:} Evaluation of DSP savings (120 DSPs vs. 181--1,512 DSPs) and BRAM reduction ($>40\times$ fewer BRAMs) enabled by on-chip streaming with zero DRAM intermediate writebacks.
    \item \textbf{Power and Energy Efficiency:} Analysis of full-chip power (1.438~W total, 0.047~W PL dynamic logic) relative to prior accelerators (2.08~W--16.85~W) to satisfy portable bedside radiography constraints.
    \item \textbf{Application Domain Focus:} Contrast between prior natural image benchmarks (Set5, Set14) and our specialized clinical radiography evaluation.
\end{itemize}

% =========================================================================
% IV-F. DISCUSSION OF PRACTICAL LIMITATIONS (OUTLINE)
% =========================================================================
\subsection{Discussion of Practical Limitations and Future Work}
\label{sec:limitations}

\begin{itemize}
    \item \textbf{Host Memory Bandwidth Bottlenecks:} Discussion of CPU-bound array slicing during canvas reassembly on large radiographs ($>1.5\text{ Mpixels}$), highlighting the future trajectory of migrating bicubic interpolation and seam exclusion into autonomous PL hardware IPs.
    \item \textbf{Fixed Magnification Topology:} Evaluation of current $2\times$ scale factor constraints and pathways to support multi-scale or reconfigurable scaling ($3\times, 4\times$).
    \item \textbf{Extreme Noise Resilience:} Considerations regarding severe Poisson noise in low-dose radiographic protocols and joint denoising-SR architectures.
\end{itemize}

\bibliographystyle{IEEEtran}
\bibliography{bib/references}
\end{document}
